% ECONOMICS_PROBLEM: {"id": "L94-166", "title": "Electricity reliability with renewables", "jel_code": "L94", "source_id": "OP-0437"}
% !TeX program = xelatex
\documentclass[11pt,a4paper]{article}
\usepackage[margin=23mm]{geometry}
\usepackage{fontspec}
\setmainfont{Latin Modern Roman}
\usepackage{amsmath,amssymb,mathtools}
\usepackage{xcolor,enumitem,booktabs,longtable,array,xurl,hyperref}
\definecolor{muted}{HTML}{53636C}
\hypersetup{colorlinks=true,urlcolor=blue,linkcolor=blue}
\setlength{\parindent}{0pt}
\setlength{\parskip}{5pt}
\newcommand{\R}{\mathbb R}
\newcommand{\E}{\mathbb E}
\newcommand{\N}{\mathbb N}
\newcommand{\ind}{\mathbf 1}
\newcommand{\Var}{\operatorname{Var}}
\newcommand{\Cov}{\operatorname{Cov}}
\newcommand{\supp}{\operatorname{supp}}
\DeclareMathOperator*{\argmin}{arg\,min}
\DeclareMathOperator*{\argmax}{arg\,max}
\newcommand{\entrymeta}[3]{{\sffamily\small\color{muted}\textbf{\texttt{#1}}\quad|\quad #2\par #3\par}}
\newcommand{\Problemref}[1]{\texttt{#1}}
\newcommand{\JELref}[2]{\texttt{#2} (JEL \texttt{#1})}
\begin{document}
\section*{Electricity reliability with renewables}
\textbf{Problem ID:} \texttt{L94-166}\quad \textbf{JEL:} \texttt{L94}\par
% BEGIN ECONOMICS STATEMENT
\label{problem:OP-0437}
\label{audited:ECON-069}
{\sffamily\small\color{muted}Original subject group: Climate, environment and energy\par}
\label{definitions:problem:30007719}
\textbf{Audit classification:} Published research agenda. Checked source identity and appropriate agenda/background support; expanded intervention, units, population, definedness and causal restrictions. Exact targets and variants are editorial.
\subsection*{Definitions and precise research target}
\textit{Editorial formalization.} Consider one interconnected electricity system with hourly periods over a representative weather year and an investment horizon of twenty years. Let \(d\in\mathcal D\) index a finite menu of implementable market designs, including a stated energy-only benchmark and specified capacity or reliability obligations. Under design \(d\), generators choose capacity and consumers may shift demand; stochastic weather, outages, and fuel prices have a prespecified joint distribution. Define total expected resource cost \(C(d)\), annual unserved energy \(L(d)\), and generation emissions \(G(d)\). For reliability threshold \(\bar L\) and emissions budget \(\bar G\), ask which design satisfies
\[d^*\in\arg\min_{d\in\mathcal D} C(d)\quad\text{subject to}\quad E[L(d)]\leq\bar L,\ E[G(d)]\leq\bar G.\]
Require finite expected costs and report infeasibility if no design meets both constraints. All costs include generation, storage, network expansion, demand response, and the specified value of lost load. Equilibrium investment and dispatch must respect transmission and storage constraints. Establish which comparisons are identified from market evidence and which rely on behavioral or technological assumptions, especially during correlated renewable shortfalls. Test design rankings under alternative scarcity-pricing limits and entry responses. This is an operational policy comparison adapted from an electricity-market research agenda, not a theorem that existing grids cannot remain reliable with renewable generation.

Let \(\omega\) collect weather, outages, and demand shocks, with declared law \(P\). For design \(d\), a measurable equilibrium-selection rule determines capacities, dispatch, load, and storage state. In hour \(t\), unserved energy is \((\operatorname{load}_t-\operatorname{served}_t)_+\) MWh; \(L(d,\omega)\) sums this over a declared representative year and \(G(d,\omega)\) sums CO\(_2\)-equivalent tonnes. \(C(d)=E_P[C_{\rm resource}(d,\omega)]\) includes a stated lost-load valuation and a common twenty-year discount convention. \(\bar L\) and \(\bar G\) refer to the same year/scenario convention. Storage obeys nonnegative state and charge/discharge capacity bounds; nodal flows obey declared transmission limits. Finite nonempty feasible menus attain their minima; identifying equilibria and rankings is the research question. All expectations use a stated baseline population and probability law. Identification means every admissible data-generating model with the same observed-data law yields the same displayed target; otherwise report the identified set. Exact equations, horizons and assumptions are editorial, rather than source-given canonical conjectures.
\subsection*{Variants and assumption changes}
\begin{enumerate}
\item \textbf{Editorial tail-reliability variant.} Replace expected unserved energy by \(P(L>\ell_0)\leq\alpha\), with \(\ell_0\geq0,\alpha\in(0,1)\). Compare design rankings at the same emissions cap.
\item \textbf{Editorial climate-trend variant.} Replace one stationary year by laws \(P_t\) allowing multiday renewable shortfalls and correlated outages. Recompute the twenty-year comparison.
\end{enumerate}
\subsection*{References and source audit}
\textbf{1.} Explicit agenda on forward obligations and real-time demand incentives; source does not solve proposed constrained designs. Locator: Report 24-09, June 2024; Section 5, printed pp.38–40. Full text or primary publication page inspected. URL: \url{https://www.rff.org/documents/4511/RFF_Report_24-09.pdf}.
\begin{enumerate}\raggedright
\item\hypertarget{definitions:source:30007719:1}{} Chiara Lo Prete; Karen Palmer; Molly Robertson. 2024. \emph{Time for a Market Upgrade? A Review of Wholesale Electricity Market Designs for the Future}. Report 24-09, June 2024; Section 5, printed pp.38–40.
\url{https://www.rff.org/documents/4511/RFF_Report_24-09.pdf}.
\end{enumerate}

\subsection*{Integrated refinement: ECON-069}
{\sffamily\small\color{muted}Renewable electricity, reliability, and productive use\par}
This refinement adopts the additional source conventions
(page~\pageref{audited:conventions}), subject to its local restrictions.
\textbf{Reliable renewable capacity and productive-use welfare.}
For feasible renewable, storage and backup capacities $k=(k_R,k_S,k_B)$ and
nonanticipative dispatch, let $B(k)$ be monetary gross final benefits,
$C(k)$ system and private adaptation resource costs, and $D(k)$ unpriced emissions
damage. Characterize $\sup_{k\in\mathcal K}\{\mathbb E[B(k)]-C(k)-D(k)\}$ and
the reliability/productive-use frontier. Specify network and storage physics,
correlated weather and outages, pricing, endogenous household and firm demand,
private backup and capital recovery. Installed capacity, dispatched renewable
share and delivered reliability differ. Internal electricity payments are
transfers; do not subtract system costs from profits already net of electricity
bills without including utility receipts. Alternatively use complete household
money-metric welfare and subtract only omitted costs/damages. Carbon payments
already included are not a second resource cost.
\subsection*{Supplied source audit for ECON-069}
Local [S1], [S2], etc. in this audit and the references immediately below
belong to ECON-069, independently of the original entry's reference numbering.
[S1] and [S2] support renewables, reliability and productive use. The feasible-capacity and corrected accounting formulation is editorial.
\subsection*{References for integrated refinement ECON-069}
{\footnotesize\raggedright
\par\noindent\textbf{[S1]} Robyn Meeks and Meera Mahadevan (2025). \emph{Electricity Infrastructure}. VoxDevLit 15(1), May 2025.\par \textit{Verified locator / source role:} Primary publisher page checked for review identity, authors, version and background. The specific published agenda is corroborated in [S2]; the exact formal target is editorial.\par \url{https://voxdev.org/voxdevlit/electricity-infrastructure}\par\smallskip
\par\noindent\textbf{[S2]} Oliver Hanney / VoxDev (living compilation). \emph{Evidence gaps: Key unanswered questions in development economics}. Accessed 8 October 2026. The page currently covers 20 topics; the ``16-topics URL is a historical slug.\par \textit{Verified locator / source role:} Topic heading: Electricity Infrastructure. Supports the broad research agenda; this is not a verbatim quotation or an attribution of the displayed mathematical target.\par \url{https://voxdev.org/topic/evidence-gaps-key-unanswered-questions-across-16-topics-development-economics}\par\smallskip
}

\subsection*{Common conventions: Source mathematical conventions}

These conventions apply to each entry unless explicitly replaced.

\textbf{Models and identification.} A primitive model \(m\) specifies all probability laws, feasible choices, information, equilibrium and measurement rules used by its target. The admissible class \(\mathcal M\) is fixed independently of outcomes. If \(P_O(m)\) is its observable probability law and \(\tau(m)\) the target, its identified set at observable law \(P\) is
\[
 \mathcal I_\tau(P)=\{\tau(m):m\in\mathcal M,\ P_O(m)=P\}.
\]
Point identification means that this set is a singleton. An empty set signals incompatibility of data and assumptions. Coordinate bounds need not describe the joint set. Bounds are sharp only if every retained target is attainable or arbitrarily approachable by compatible models. Finite bounds require explicit restrictions on unobserved quantities; unrestricted confounding, hidden wealth, extrapolation or arbitrary equilibrium selection can yield uninformative sets.

\textbf{Causal interventions.} A potential outcome \(Y_i(p)\) is the outcome under a complete policy regime \(p\), including timing, financing, information and market spillovers. The target population and units are fixed before treatment. With interference, use the complete assignment vector or a justified exposure mapping; individual no-interference notation alone cannot identify an economy-wide intervention. A difference of mean potential outcomes does not require the cross-world joint distribution. Conditioning on post-treatment migration, survival, enrollment or employment changes the estimand and needs separate justification.

\textbf{Statistical statements.} An estimator, confidence set, forecast or test specifies a sampling law, model class and loss/coverage criterion. Uniform \((1-\alpha)\)-coverage means \(\inf_{m\in\mathcal M}\Pr_m\{\tau(m)\in C_n\}\geq1-\alpha\), or an explicitly asymptotic version. The whole parameter space trivially has coverage, so informativeness requires a stated width, risk or power objective. A forecasting improvement is relative to a named benchmark, information set and evaluation population.

\textbf{Equilibrium and optimization.} An equilibrium comprises feasible plans, optimality, beliefs where needed, budget constraints and all market/resource clearing. The primitives or a stated benchmark define each correspondence \(\mathcal E(p;m)\). Nonemptiness is assumed or proved, never inferred just from compact policy sets. A selected-equilibrium target uses a declared selection rule; a robust target quantifies over all permitted equilibria. Use suprema or infima unless attainment is established. An \(\varepsilon\)-optimal policy is feasible and within \(\varepsilon\) of the finite optimum value. Report an empty feasible set separately.

\textbf{Welfare and accounting.} Normative utility, interpersonal normalization, social weights, numeraire, discounting and the use of public receipts are primitives. Behavioral utility need not coincide with the welfare criterion. Household welfare includes ownership income and rents once; adding profits again to compensating variation generally double-counts them. A monetary equivalent is defined by an explicit transfer and price convention, with monotonicity and range conditions. Average effects, mechanism contributions, transfers and welfare losses are different targets. Infinite sums require convergence and dynamic feasible plans require terminal or no-Ponzi conditions.

\textbf{Source versions.} A working-paper date, revision date and journal publication year are distinguished. A DOI for a journal version is not automatically the DOI of an earlier manuscript. Locators refer to the stated version and printed pages where specified. A metadata-only check does not verify a claim about a full-text conclusion. Every entry's source subsection narrows the provenance claim accordingly.

\subsection*{Common conventions: Additional-source status and mathematical conventions}

\label{audited:conventions}

Of the 100 supplied ECON entries, 65 distinct problems appear here; 32 close
matches refine earlier entries, and three duplicate questions map to existing
entries. Appendix~\ref{app:audited-crosswalk} lists every destination.
The attachment's P/R/S/U distinctions and source audits are retained. Its
precise mathematical formalizations are editorial unless the individual audit
says otherwise. Candidate subsidy targets ECON-004--006 retain their U flags;
the resolved approval-core question ECON-007 updates OP-0202 above.
The following supplied conventions govern both this part and the attributed
ECON refinements elsewhere in the catalogue, subject to local restrictions.

\textbf{Parameterized questions.} A problem family lists inputs that must be fixed before an instance is evaluated: population, horizon, policies, information, data law, model class and welfare weights. These are required choices, not quantities the citations are assumed to specify. Undefined applications should not be treated as identified estimands.

\textbf{Indices and integrability.} Sets of agents, firms and regions are finite unless a population measure is explicitly used. \(i,j\) index units, \(r\) regions, \(t\) dates and \(h\) horizons. \(\mathbb E\) and \(\Pr\) refer to the declared population and law; \(\mathbf1\{A\}\) is an indicator. Every displayed expectation and discounted sum needs existence in the stated domain. Infinite horizons use a discount in \((0,1)\) and a convergence condition; finite horizons may use discount one.

\textbf{Optimization and existence.} A displayed \(\max\), \(\min\), or \(\arg\max\) is conditional on attainment. For finite feasible menus this is immediate; general spaces need continuity and compactness/coercivity or another existence argument. Without attainment, the target is the supremum/infimum and arbitrarily near-optimal feasible choices. \(\inf\varnothing=+\infty\). Empty feasibility and an unidentified target are different issues.

\textbf{Potential outcomes.} \(Y_i(z)\) is the outcome under a specified intervention, not the observed conditional mean among people choosing \(z\). In binary comparisons \(\Delta Y_i=Y_i(1)-Y_i(0)\). Specify assignment, treatment versions, compliance, information and observation horizon. Interference is allowed under a full policy vector; compressed exposure notation needs a justified mapping. No interference, consistency and overlap are substantive assumptions, not automatic consequences of notation.

\textbf{Populations and observation.} Fix baseline eligibility and population weights. Conditioning on post-policy employment, survival, relocation or program uptake can change the population being compared. Death is not ordinary loss to follow-up; outcomes truncated by death need a composite convention, joint survival target, or explicitly defined survivor stratum. Fertility-changing interventions need a population functional rather than an unexplained fixed descendant cohort. Measurement and missingness models must be supplied.

\textbf{Identification.} A structural model \(m\in\mathcal M\) implies observable law \(P_m\) and target \(\theta(m)\). Identification means
\[
 P_{m_1}=P_{m_2}\ \Longrightarrow\ \theta(m_1)=\theta(m_2)
 \quad\text{for all }m_1,m_2\in\mathcal M .
\]
Without identification, the target is
\[
 \Theta(P;\mathcal M)=\{\theta(m):m\in\mathcal M,\ P_m=P\}.
\]
Writing \(\theta(P)\) before establishing identification can hide incompatible latent models. Confidence sets additionally need a sampling law, confidence level, asymptotic sequence and uniformity class. Point identification, coverage and informative precision are separate properties.

\textbf{Mediation.} Total effects of randomized assignment are distinct from cross-world natural indirect effects. Belief/effort-channel effects require a meaningful mediator intervention and additional measurement, exclusion or mediation assumptions. Randomization of the initial policy alone does not supply those assumptions. Report bounds or interventional alternatives when the stronger target is unsupported.

\textbf{Equilibrium.} A counterfactual \(W(z)\), \(p(z)\), or \(\mu_t^z\) requires individual optimization, feasibility, government/resource budgets, market clearing and state-transition laws. State equilibrium existence and selection, or report ranges over compatible equilibria. Initial-state integration includes subsequent endogenous transitions; current-state integration must use the policy-dependent distribution. Reduced-form invariance after policy is not assumed.

\textbf{Welfare accounts.} Scores, utility, health, dollars and ecological quantities require explicit conversion before addition. Fees, taxes, subsidies, wages and extortion payments are transfers between parties; distributional weights can make them welfare relevant, but their face value is not automatically a resource loss. State ownership, geographic boundaries, foreign-income weights and fiscal finance. Do not add profits to household welfare already including dividends, or count land capitalization and the underlying benefit twice. A benefit-cost expression must distinguish gross benefits from net welfare and subtract every real cost exactly once.

\textbf{Derivatives and risk.} Log derivatives require positive arguments; local responses require admissible perturbations and specified equilibrium branches. Differentiation under expectations needs regularity/domination, and boundaries may allow only one-sided derivatives. Risk uses a specified law; ambiguity uses a declared nonempty law set on a common history space. Adaptive policies must be nonanticipative. A robust criterion is an editorial choice unless attributed explicitly.

\textbf{Scope overlaps.} Allocation, collective choice, contracts, household bargaining and coordinated policy overlap economics with optimization and game theory. They are retained because they appeared in the original catalogue; no claim of a strictly game-theory-free selection is made.
% END ECONOMICS STATEMENT
\end{document}
